%% ============================================================
%% Chapter 5 -- Experimental Results and Analysis
%% ============================================================

\chapter{Experimental Results and Analysis}\label{chapter:5_Experiments}

This chapter presents the experimental results for the six tone-classification
setups defined in Chapter~\ref{chapter:4_Design}. The previous chapter described
the CNN baseline, the Wav2Vec~2.0 adaptation routes, the HuBERT variants, and
the speaker-disjoint learner split used by the self-supervised systems. This
chapter evaluates how well those systems reproduce expert perceptual tone labels
on held-out Polish learner speech. The analysis first defines the evaluation
metrics and protocol, then compares the aggregate results, examines training
behaviour, analyses the most difficult tone categories, and finally discusses the
findings with respect to the research questions stated in
Chapter~\ref{chapter:3_Proposal}.

\section{Evaluation Metrics and Protocol}
\label{sec:exp:protocol}

The evaluation measures agreement between the predicted tone category
\(\hat{y}_i\) and the expert label \(y_i\) for each test utterance \(i\). The
target set is \(\mathcal{T}=\{T_1,T_2,T_3,T_4,T_{\mathrm{NS}}\}\), where
\(T_{\mathrm{NS}}\) denotes a non-standard learner production rather than the
linguistic neutral tone. Accuracy reports the proportion of correctly classified
utterances:
\begin{equation}
  \mathrm{Accuracy}
  =
  \frac{1}{N}\sum_{i=1}^{N}\mathbf{1}\!\left[\hat{y}_i=y_i\right].
  \label{eq:exp:accuracy}
\end{equation}

Because the learner corpus is class-imbalanced, accuracy is not sufficient on
its own. For each class \(c\), precision, recall, and F1 are computed as
\begin{equation}
  P_c=\frac{\mathrm{TP}_c}{\mathrm{TP}_c+\mathrm{FP}_c},
  \qquad
  R_c=\frac{\mathrm{TP}_c}{\mathrm{TP}_c+\mathrm{FN}_c},
  \qquad
  F1_c=\frac{2P_cR_c}{P_c+R_c}.
  \label{eq:exp:class_f1}
\end{equation}
Macro-F1 gives equal weight to all five tone categories, while weighted F1
weights each class by its support \(n_c\):
\begin{equation}
  \mathrm{Macro}\text{-}F1
  =
  \frac{1}{|\mathcal{T}|}\sum_{c\in\mathcal{T}}F1_c,
  \qquad
  \mathrm{Weighted}\ F1
  =
  \sum_{c\in\mathcal{T}}\frac{n_c}{N}F1_c.
  \label{eq:exp:aggregate_f1}
\end{equation}
Macro-F1 is therefore the most informative aggregate metric for this thesis,
because it exposes whether a model performs poorly on sparse categories such as
\(T_3\) or \(T_{\mathrm{NS}}\).

Model selection always used development data, and the final test set was
evaluated only after checkpoint selection was complete. The CNN baseline uses
a random image split with $N=1{,}194$ test images, selecting its checkpoint by
validation accuracy. The five SSL setups share a speaker-disjoint learner
split with $N=2{,}282$ test utterances and 464 development utterances; their
selection metric is learner-development accuracy for Setups~2, 3a--3c (learner
stage), 4, and 5, and learner-development macro-F1 for Setup~6. For the
intermediate stage in Setups~3a--3c the Stage~A checkpoint is picked on the
AISHELL-1 development set (693 utterances, frame accuracy) before the learner
stage takes over.

\section{Overall Model Comparison}
\label{sec:exp:aggregate}

Table~\ref{tab:exp:aggregate} reports the aggregate test results. The main
comparison is among the five self-supervised systems because they use the same
speaker-disjoint learner test split. The CNN baseline is included to show the
performance of a conventional mel-spectrogram classifier under its original
random image split.

\begin{table}[H]
\centering
\caption{Aggregate test performance for the six experimental setups. SSL rows
share the speaker-disjoint learner test set ($N=2\,282$); the CNN row uses an
independent random image split ($N=1\,194$) and is reported as a ToneNet-style
reference. Higher values are better; best values per metric are in
\textbf{bold}.}
\label{tab:exp:aggregate}
\scriptsize
\setlength{\tabcolsep}{4pt}
\begin{tabular}{@{}>{\centering\arraybackslash}p{0.05\textwidth}
>{\raggedright\arraybackslash}p{0.36\textwidth}
>{\raggedright\arraybackslash}p{0.20\textwidth}
>{\centering\arraybackslash}p{0.08\textwidth}
>{\centering\arraybackslash}p{0.08\textwidth}
>{\centering\arraybackslash}p{0.08\textwidth}@{}}
\toprule
\textbf{Setup} & \textbf{System} & \textbf{Pre-training source} & \textbf{Acc. (\%)} & \textbf{Macro-F1 (\%)} & \textbf{Weighted F1 (\%)} \\
\midrule
1  & ToneNet-style CNN                                & n/a                          & 72.53 & 60.00 & 71.00 \\
2  & Wav2Vec~2.0, learner fine-tuning                 & LibriSpeech                  & 77.56 & 70.64 & 77.39 \\
3a & Wav2Vec~2.0 + AISHELL-1 (label1)                 & LibriSpeech $+$ AISHELL-1    & 77.13 & 71.46 & 77.18 \\
3b & Wav2Vec~2.0 + AISHELL-1 (label2)                 & LibriSpeech $+$ AISHELL-1    & 76.69 & 71.41 & 76.65 \\
3c & Wav2Vec~2.0 + AISHELL-1 (label3)                 & LibriSpeech $+$ AISHELL-1    & 76.91 & 71.42 & 77.01 \\
4  & Chinese Wav2Vec~2.0-Large                        & WenetSpeech       & 77.26 & 67.33 & 75.49 \\
5  & Chinese HuBERT-Large, frame-level CE             & WenetSpeech       & \textbf{80.28} & \textbf{73.37} & \textbf{80.00} \\
6  & Chinese HuBERT-Large, attention $+$ focal loss   & WenetSpeech       & 78.62 & 71.61 & 78.40 \\
\bottomrule
\end{tabular}
\end{table}

The strongest overall system is Setup~5, the Chinese HuBERT-Large model trained
with frame-level cross-entropy and majority voting. It reaches 80.28\,\%
accuracy, 73.37\,\% macro-F1, and 80.00\,\% weighted F1. Compared with the
direct English-pretrained Wav2Vec~2.0 learner fine-tuning baseline (Setup~2),
this corresponds to gains of 2.72 percentage points in accuracy and 2.73
percentage points in macro-F1. Setup~6 uses the same HuBERT backbone but replaces
frame voting with attention pooling and focal loss. It still performs strongly,
with 78.62\,\% accuracy and 71.61\,\% macro-F1, but it does not improve over the
simpler frame-level HuBERT objective.

The Wav2Vec~2.0 results show that intermediate native-Mandarin adaptation does
provide a modest class-balanced benefit, but not a clear accuracy gain. Direct
learner fine-tuning (Setup~2) obtains 77.56\,\% accuracy and 70.64\,\%
macro-F1. The three AISHELL-1 streams converge to a narrow band of
71.41--71.46\,\% macro-F1, all above the learner-only baseline, while their
accuracies remain slightly below Setup~2. The label1 stream (Setup~3a) reaches
the strongest downstream result among the three AISHELL-1 variants, with
77.13\,\% accuracy and 71.46\,\% macro-F1. The class-level analysis in
Section~\ref{sec:exp:class_level} shows that the AISHELL-1 intermediate stage
improves both \(T_3\) and \(T_{\mathrm{NS}}\) behaviour, while modest losses on
some canonical tones limit its accuracy gains.

The placement of frame labels within the AISHELL-1 syllable (label1, label2,
label3) has a smaller effect on downstream learner accuracy than the presence
of the AISHELL-1 stage itself. The three streams reach Stage~A frame accuracy
of 87.90\,\%, 92.05\,\%, and 91.00\,\% on the AISHELL-1 development set;
label2 is the highest of the three, but label1 obtains the strongest
downstream learner result. The native-frame selection score therefore does not
perfectly predict learner-test performance, even though all three streams remain
close to one another after learner fine-tuning. On the learner development set,
the SSL backbones reach 75.65\,\% (Setup~2), 76.29\,\% (Setup~4), and 81.90\,\%
(Setup~5) accuracy and 73.86\,\% macro-F1 for the attention head (Setup~6);
these dev scores match the ordering of the held-out test results, indicating
that checkpoint selection was stable.

The Mandarin-pretrained Wav2Vec~2.0 checkpoint (Setup~4) reaches 77.26\,\%
accuracy, comparable to the AISHELL-1 routes, but its macro-F1 falls to
67.33\,\%, the lowest among the SSL systems. The drop is caused by very poor
recognition of \(T_{\mathrm{NS}}\), as shown in
Section~\ref{sec:exp:class_level}. Therefore, language-matched pre-training
alone does not guarantee class-balanced improvement; the downstream architecture
and objective also determine whether the model can handle non-canonical learner
tones.

\section{Training Behaviour}
\label{sec:exp:curves}

The training curves were examined to verify that the selected checkpoints did
not come from isolated spikes in development performance.

The CNN baseline overfits its training set
(Figure~\ref{fig:exp:cnn_curve}). Training accuracy reaches 100\,\% within
roughly ten epochs while validation accuracy plateaus near 68\,\% and
validation loss starts to creep up. This is the textbook overfit signature for
a from-scratch image classifier on a small split, and explains why the CNN
checkpoint is selected on validation accuracy rather than left at the final
epoch. The confusion matrix on the same run
(Figure~\ref{fig:exp:cnn_cm}) shows that the residual error is concentrated on
\(T_3\) and \(T_{\mathrm{NS}}\): \(T_3\) is heavily confused with \(T_2\), and
\(T_{\mathrm{NS}}\) productions are scattered across all four canonical tones.
This per-class pattern is the same one quantified for the CNN row of
Tables~\ref{tab:exp:aggregate} and~\ref{tab:exp:class_level}, although the
absolute counts there come from an earlier CNN run on the smaller
$N=1{,}194$ image split rather than the run plotted here.

\begin{figure}[H]
\centering
\safeincludegraphics[width=0.85\textwidth]{code/output/cnn_training/cnn_training_curve.png}
\caption{CNN baseline training and validation curves.}
\label{fig:exp:cnn_curve}
\end{figure}

\begin{figure}[H]
\centering
\safeincludegraphics[width=0.6\textwidth]{code/output/cnn_training/cnn_training_confusion_matrix.png}
\caption{CNN baseline test-set confusion matrix.}
\label{fig:exp:cnn_cm}
\end{figure}

The SSL curves
(Figures~\ref{fig:exp:w2v_direct}, \ref{fig:exp:aishell_curves},
\ref{fig:exp:aishell_learner_curves}, and~\ref{fig:exp:hubert_curves}) all use
a vertical dashed line to mark the point at which the Transformer blocks were
unfrozen after the initial head-only training phase

\begin{figure}[H]
\centering
\begin{minipage}{0.49\textwidth}
\centering
\safeincludegraphics[width=\linewidth]{code/output/w2v2_large_en_baseline_learner/training_curves_smoothed.png}
\small English checkpoint Wav2Vec~2.0
\end{minipage}
\hfill
\begin{minipage}{0.49\textwidth}
\centering
\safeincludegraphics[width=\linewidth]{code/output/w2v2_large_zh_learner/training_curves_smoothed.png}
\small Chinese checkpoint Wav2Vec~2.0-Large
\end{minipage}
\caption{Training curves for direct learner fine-tuning of Wav2Vec~2.0 systems.}
\label{fig:exp:w2v_direct}
\end{figure}

\begin{figure}[H]
\centering
\safeincludegraphics[width=\textwidth]{code/output/wave2vec_aishell1_training_curves_smoothed.png}
\caption{AISHELL-1 intermediate-training curves for the three frame-label
streams used in Setups~3a--3c.}
\label{fig:exp:aishell_curves}
\end{figure}

\begin{figure}[H]
\centering
\safeincludegraphics[width=\textwidth]{code/output/wave2vec_aishell1_learner_dataset_training_curves_smoothed.png}
\caption{Learner fine-tuning curves for Wav2Vec~2.0 checkpoints initialised
from AISHELL-1.}
\label{fig:exp:aishell_learner_curves}
\end{figure}

\begin{figure}[H]
\centering
\begin{minipage}{0.49\textwidth}
\centering
\safeincludegraphics[width=\linewidth]{code/output/hubert_large_zh_learner/training_curves_smoothed.png}
\small Frame-level cross-entropy
\end{minipage}
\hfill
\begin{minipage}{0.49\textwidth}
\centering
\safeincludegraphics[width=\linewidth]{code/output/hubert_large_zh_attn_learner/training_curves_smoothed.png}
\small Attention pooling with focal loss
\end{minipage}
\caption{Training curves for Chinese HuBERT-Large learner fine-tuning.}
\label{fig:exp:hubert_curves}
\end{figure}

The Wav2Vec~2.0 curves show the expected two-phase pattern. During the frozen
Transformer stage, the classifier head learns a stable initial mapping from SSL
features to tone categories. After unfreezing, development accuracy rises and
then levels off, indicating that the selected checkpoints are part of a stable
training region rather than one-off fluctuations. The AISHELL-1 curves reach
higher frame-level development accuracy than the learner curves, which is
expected because the intermediate task uses clean native speech and frame-level
labels. However, the learner results in Table~\ref{tab:exp:aggregate} show that
the highest AISHELL-1 development score does not necessarily produce the best
learner classifier.

Both HuBERT variants train smoothly after unfreezing, but their final test
results differ. The frame-level HuBERT model achieves the best development
accuracy and also the strongest held-out test performance. The attention-pooling
model is selected by macro-F1 and is designed to reduce the effect of class
imbalance, yet its held-out macro-F1 remains below the frame-level model. In
this dataset, the additional pooling mechanism did not compensate for the
benefit of dense frame-level supervision and majority voting.

\section{Class-Level Error Analysis}
\label{sec:exp:class_level}

Aggregate scores hide important differences between tone categories. Table
\ref{tab:exp:class_level} reports per-class recall and per-class \(F_1\) for
every class and every setup. It is the consolidated class-level reference for
the rest of this chapter. The two categories \(T_{\mathrm{NS}}\) and
\(T_3\) produced the largest class-level errors across the SSL systems and are
also linguistically important: \(T_{\mathrm{NS}}\) groups non-canonical
learner productions, and \(T_3\) is the canonical tone most often
mis-produced by learners whose first language does not use lexical
tone~\cite{hao2012}.

\begin{table}[H]
\centering
\caption{Per-class precision (P), recall (R), and \(F_1\) on the
speaker-disjoint learner test split (CNN uses the random image split). Values
are in percentage points; the strongest value for each class--metric pair
(excluding the CNN reference) is in \textbf{bold}. CNN precision is recovered
from the reported R and \(F_1\) because the original CNN run is not retained in
machine-readable form. Confusion matrices for every setup are reproduced in
Appendix~\ref{sec:app:confusion_matrices}.}
\label{tab:exp:class_level}
\tiny
\setlength{\tabcolsep}{2pt}
\resizebox{\textwidth}{!}{%
\begin{tabular}{@{}>{\centering\arraybackslash}p{0.03\textwidth}
>{\raggedright\arraybackslash}p{0.21\textwidth}
ccccccccccccccc@{}}
\toprule
 & & \multicolumn{3}{c}{\(T_{\mathrm{NS}}\)} & \multicolumn{3}{c}{\(T_1\)} & \multicolumn{3}{c}{\(T_2\)} & \multicolumn{3}{c}{\(T_3\)} & \multicolumn{3}{c}{\(T_4\)} \\
\cmidrule(lr){3-5}\cmidrule(lr){6-8}\cmidrule(lr){9-11}\cmidrule(lr){12-14}\cmidrule(lr){15-17}
\textbf{S} & \textbf{System} & \textbf{P} & \textbf{R} & \(F_1\) & \textbf{P} & \textbf{R} & \(F_1\) & \textbf{P} & \textbf{R} & \(F_1\) & \textbf{P} & \textbf{R} & \(F_1\) & \textbf{P} & \textbf{R} & \(F_1\) \\
\midrule
1  & CNN baseline                       & 60.79 & 37.00 & 46.00 & 76.03 & 78.00 & 77.00 & 72.57 & 82.00 & 77.00 & 51.33 & 14.00 & 22.00 & 71.36 & 86.00 & 78.00 \\
2  & Wav2Vec~2.0, learner only          & 55.92 & 47.58 & 51.42 & 80.36 & 80.60 & 80.48 & 81.91 & 79.78 & 80.83 & 50.00 & 57.60 & 53.53 & 84.13 & 89.94 & 86.94 \\
3a & Wav2Vec~2.0 + AISHELL-1 label1     & 59.91 & 53.63 & 56.60 & 79.50 & 80.45 & 79.97 & 81.63 & 76.50 & 78.98 & 48.50 & 64.80 & 55.48 & 84.50 & 88.17 & 86.29 \\
3b & Wav2Vec~2.0 + AISHELL-1 label2     & 61.97 & 53.23 & 57.27 & 77.30 & 80.30 & 78.77 & 81.97 & 75.14 & 78.40 & 50.31 & \textbf{65.60} & \textbf{56.94} & 83.12 & 88.36 & 85.66 \\
3c & Wav2Vec~2.0 + AISHELL-1 label3     & 57.08 & \textbf{55.24} & 56.15 & 79.06 & 80.60 & 79.82 & 80.61 & 76.09 & 78.29 & 50.32 & 63.20 & 56.03 & 86.50 & 87.18 & 86.84 \\
4  & Chinese Wav2Vec~2.0-Large          & \textbf{72.09} & 25.00 & 37.13 & 79.46 & 83.13 & 81.25 & 74.40 & \textbf{84.97} & 79.34 & \textbf{59.18} & 46.40 & 52.02 & 82.71 & \textbf{91.52} & 86.89 \\
5  & Chinese HuBERT-Large, CE           & 66.32 & 50.81 & \textbf{57.53} & \textbf{81.78} & \textbf{85.07} & \textbf{83.39} & \textbf{83.02} & 83.47 & \textbf{83.24} & 52.59 & 56.80 & 54.62 & \textbf{86.64} & 89.55 & \textbf{88.07} \\
6  & Chinese HuBERT-Large, attn $+$ focal & 64.97 & 51.61 & \textbf{57.53} & 79.89 & 84.18 & 81.98 & 81.50 & 80.05 & 80.77 & 48.18 & 52.80 & 50.38 & 86.04 & 88.76 & 87.38 \\
\bottomrule
\end{tabular}%
}
\end{table}

The clearest class-level failure appears in Setup~4. Chinese Wav2Vec~2.0-Large
correctly identifies only 62 of the 248 \(T_{\mathrm{NS}}\) test utterances,
giving 25.00\,\% recall. The confusion matrix
(Table~\ref{tab:exp:cm_w2v_chinese}) shows that 91 of those 248 utterances are
assigned to \(T_2\). The same backbone produces the strongest \(T_2\) recall
of any setup (84.97\,\%), confirming that the representation is rich for
canonical Mandarin tones; the deficit is in modelling the absence of canonical
tone. This pattern is consistent with a Mandarin-pretrained backbone whose
internal templates for the four canonical tones are sharp enough to absorb
acoustically ambiguous learner productions into the nearest canonical class,
rather than preserving the expert's non-standard judgement.

\begin{table}[H]
\centering
\caption{Setup~4 (Chinese Wav2Vec~2.0-Large) confusion matrix on the
speaker-disjoint learner test split. Rows are expert labels; columns are
model predictions.}
\label{tab:exp:cm_w2v_chinese}
\small
\begin{tabular}{@{}lrrrrr@{}}
\toprule
 & \textbf{Pred. \(T_{\mathrm{NS}}\)} & \textbf{Pred. \(T_1\)} & \textbf{Pred. \(T_2\)} & \textbf{Pred. \(T_3\)} & \textbf{Pred. \(T_4\)} \\
\midrule
True \(T_{\mathrm{NS}}\) & 62 & 27 & 91 & 25 & 43 \\
True \(T_1\)             &  9 & 557 & 73 &  0 & 31 \\
True \(T_2\)             &  8 & 82 & 622 & 10 & 10 \\
True \(T_3\)             &  5 & 11 & 38 & 58 & 13 \\
True \(T_4\)             &  2 & 24 & 12 &  5 & 464 \\
\bottomrule
\end{tabular}
\end{table}

For \(T_3\), all three AISHELL-1-initialised Wav2Vec~2.0 systems improve recall
over the learner-only baseline (Setup~2): 64.80, 65.60, and 63.20\,\% against
57.60\,\%. Setup~3b obtains the best \(T_3\) recall and \(F_1\) (56.94\,\%)
among all SSL systems, while Setups~3a and 3c stay close behind. Native Mandarin
frame supervision therefore selectively helps the canonical tone most often
mis-produced by non-tonal-L1 learners.

For \(T_{\mathrm{NS}}\), the AISHELL-1 routes and the HuBERT routes both improve
over the learner-only Wav2Vec~2.0 baseline. Setup~3c gives the highest
\(T_{\mathrm{NS}}\) recall (55.24\,\%), while Setups~5 and~6 tie for the highest
\(T_{\mathrm{NS}}\) \(F_1\) (57.53\,\%). The attention-focal HuBERT variant is
therefore not uniquely strongest on this class, and its lower \(T_3\) score
helps explain why it remains below the simpler HuBERT CE model on aggregate
metrics. The two minority categories are not improved jointly by any single
setup in this study. The remaining \(T_3\) and \(T_{\mathrm{NS}}\) \(F_1\) values,
all in the 50--58\,\% range, are well
below the \(F_1\) values of \(T_1\), \(T_2\), and \(T_4\) (typically above
80\,\% across the SSL systems), and these two categories therefore remain
the main bottlenecks for an expert-aligned Mandarin tone assessment model.

\section{Discussion and Analysis}
\label{sec:exp:interpretation}

The aggregate and class-level results answer the research questions from
Chapter~\ref{chapter:3_Proposal} with different levels of certainty. The first
four questions are directly tested by the experiments, while the last two are
answered by interpreting the reported results in the broader CAPT setting.

\paragraph{RQ1: CNN versus self-supervised speech models.}
The self-supervised systems provide stronger class-balanced agreement with the
expert tone labels than the ToneNet-style CNN reference. The CNN reaches
60.00\,\% macro-F1 on its random image split, while every SSL setup reaches
at least 67.33\,\% macro-F1 on the stricter speaker-disjoint learner
protocol, and the best SSL system (Setup~5) reaches 73.37\,\%. The class-level
breakdown in Table~\ref{tab:exp:class_level} shows that the gap is
concentrated on the two minority categories: for \(T_1\) and \(T_2\) the CNN
reaches \(F_1\) values within a few percentage points of the SSL systems
(CNN 77\,\% vs SSL 78.29--83.39\,\%), and for \(T_4\) the CNN trails the SSL
systems by roughly 8 percentage points (CNN 78\,\% vs SSL 85.66--88.07\,\%),
but the largest gaps appear on \(T_3\) (CNN 22\,\% vs SSL 50.38--56.94\,\%)
and \(T_{\mathrm{NS}}\) (CNN 46\,\% vs SSL 51.42--57.53\,\%, with Setup~4 the
exception at 37.13\,\%). The class-balanced advantage of SSL therefore
concentrates on \(T_3\) and \(T_{\mathrm{NS}}\) but is also visible on \(T_4\).
Because the CNN and SSL systems do not share the same split protocol, this
result is evidence that SSL representations are more promising for this task,
not a fully controlled head-to-head comparison.

\paragraph{RQ2: AISHELL-1 intermediate adaptation.}
Intermediate adaptation on native Mandarin speech modestly lifts aggregate
macro-F1 above the learner-only baseline, although it does not improve overall
accuracy. Direct learner fine-tuning (Setup~2) reached 70.64\,\% macro-F1,
while the three AISHELL-1 variants converged to 71.41--71.46\,\% macro-F1.
The class-level analysis shows why: all three AISHELL-1 streams improved both
\(T_3\) and \(T_{\mathrm{NS}}\) recall over Setup~2, but their gains on these
minority categories were partly offset by weaker scores on some canonical
tones. Native Mandarin frame supervision therefore helps class-balanced learner
assessment, but the transfer is not strong enough to dominate the frame-level
HuBERT system. The
narrow gap among the three label streams shows that the placement of frame
labels within the syllable (full toned onset+rhyme, rhyme-only with a
separate consonant class, or midpoint-only) has a smaller effect on
downstream learner accuracy than the AISHELL-1 stage as a whole. Within the
Setup~3 group, the highest AISHELL-1 frame accuracy (92.05\,\% for label2)
does not coincide with the strongest learner-test result, which is obtained by
label1.

\paragraph{RQ3: Mandarin-specific pre-training.}
Mandarin-specific pre-training helps only when the downstream objective can also
represent non-canonical learner productions. Chinese Wav2Vec~2.0-Large reaches
accuracy close to English-pretrained Wav2Vec~2.0, but its weak
\(T_{\mathrm{NS}}\) recall lowers macro-F1. Chinese HuBERT-Large does not show
the same failure pattern and gives the best aggregate and class-level balance.
The result therefore supports Mandarin-specific pre-training for HuBERT in this
setup, but not as a universal explanation for improved learner-tone assessment.

\paragraph{RQ4: Most difficult expert categories.}
The remaining errors concentrate in \(T_3\) and \(T_{\mathrm{NS}}\). \(T_3\)
remains difficult because its contour is phonetically variable and frequently
confused by non-tonal-L1 learners. \(T_{\mathrm{NS}}\) is difficult for a
different reason: it is not a single acoustic category, but a perceptual label
for productions that do not realise any canonical tone. Future systems should
therefore target these categories directly, for example through class-aware
losses, confidence calibration, or more precise tone-bearing interval labels.

\paragraph{RQ5: Expert labels versus intended-tone evaluation.}
Evaluating against expert perceptual labels changes the interpretation of the
scores. A model is not rewarded for predicting the prompt tone if the learner's
realised production was judged differently by the expert. This is especially
important for \(T_{\mathrm{NS}}\), where the correct output is a non-standard
assessment category rather than one of the four canonical tones. The reported
accuracy and F1 values should therefore be read as agreement with expert
judgement in a CAPT context, not as ordinary native-speaker tone recognition.

\paragraph{RQ6: Integration with a phonetic-accuracy classifier.}
The tonal classifier can be combined with an independently trained
phonetic-accuracy classifier only if both modules expose reliable confidence
information and use compatible utterance or token references. The present
results show that tone prediction is strongest for canonical tones and weaker
for \(T_3\) and \(T_{\mathrm{NS}}\). A two-dimensional CAPT feedback signal
should therefore avoid overconfident corrective feedback for these difficult
classes unless it also includes calibration, prompt-tone context, and a clear
policy for disagreement between the tone and segmental classifiers.

\section{Limitations}
\label{sec:exp:limitations}

The main limitation is that the CNN baseline and the SSL systems do not share
the same split protocol. The CNN row is therefore used only as a reference point
for a conventional ToneNet-style mel-spectrogram classifier. A fully controlled
CNN-versus-SSL comparison would require the image model to be trained and tested
on the same speaker-disjoint partition used for the SSL systems.

The learner frame-level objective is also approximate. A single expert
utterance-level tone label is copied to all MFA-selected speech frames. This
provides a practical supervision signal, but it does not identify the exact
tone-bearing interval within each learner utterance. The AISHELL-1 label-stream
experiments show that temporal label placement matters, so future work should
test syllable-level or phone-level tone targets, CTC-style objectives, and
class-aware loss weighting for \(T_3\) and \(T_{\mathrm{NS}}\).

Finally, the evaluation measures agreement with one expert annotation source.
This matches the assessment target used in the thesis, but it does not estimate
inter-rater reliability. Additional expert labels would make it possible to
separate model errors from ambiguous human judgements, especially for
non-standard learner productions.

\section{Chapter Summary}
\label{sec:exp:summary}

This chapter evaluated six Mandarin tone-classification setups on Polish
learner speech. The best system was Setup~5, Chinese HuBERT-Large with
frame-level cross-entropy, which achieved 80.28\,\% accuracy, 73.37\,\%
macro-F1, and 80.00\,\% weighted F1 on the speaker-disjoint learner test set.
Among the Wav2Vec~2.0 systems, the strongest macro-F1 came from the AISHELL-1
label1 route (Setup~3a) at 71.46\,\%, with the three AISHELL-1 intermediate
variants converging to a narrow band of 71.41--71.46\,\% macro-F1; the
AISHELL-1 stage improved \(T_3\) and \(T_{\mathrm{NS}}\) behaviour but did not
improve accuracy over the learner-only baseline. The Chinese Wav2Vec~2.0
backbone (Setup~4) reached test accuracy comparable to the AISHELL-1 routes
but suffered a sharp \(T_{\mathrm{NS}}\) collapse, illustrating that the
benefit of language-matched pre-training depends on the downstream backbone.
The attention-focal HuBERT variant (Setup~6) tied Setup~5 on
\(T_{\mathrm{NS}}\) \(F_1\), but it reduced \(T_3\) and aggregate performance.
Training curves
showed stable convergence after Transformer unfreezing, and class-level
analysis identified \(T_3\) and \(T_{\mathrm{NS}}\) as the most difficult
categories. The results support self-supervised speech representations for
expert-aligned Mandarin tone assessment, but they also show that future CAPT
systems must handle non-canonical learner productions more explicitly.
